\chapter{Einleitung}
\label{ch:einleitung}

\section{Motivation und Aufgabenstellung}
Grundlage der Projektarbeit ist ein Frequenzumrichter, welcher 2023 von zwei Masterstudent der Hochschule Bochum modellbasiert entwickelt wurde. Initiale Aufgabe war es einen Anwendungsfall für den Frequenzumrichter zu realisieren.
Nach erster Bestandsaufnahme und Sichtung der Masterarbeiten stand jedoch schnell fest, dass zunächst die Firmware und Hardware überarbeitet werden muss.
Perspektivisch soll der Frequenzumrichter Firmen als Plattform für ihre eigene Entwicklung bereitgestellt werden können.

Aus dieser neu formulierten Aufgabenstellung heraus wurde das Ziel neu definiert: Ausbau des bestehenden Projekts hin zu einem FOC-Framework bestehend aus:

\begin{itemize}
	\item Wiederverwendbare C und MatLab Simulink Bibliothek für Anwendungsfälle
	\item Parametrisierung von Motoren über austauschbare Motorprofile 
	\item Ausbau der modellbasierten Tests
	\item Erstellung einer \ac{GUI}-Anwendung für die Inbetriebnahme und Bedienung des Frequenzumrichters über \ac{CAN}
\end{itemize}
\offen{Überarbeiten}

\section{Abgrenzung und Aufteilung im Team}
Dieser Projektbericht umfasst die Softwareseite der Projektarbeit.
Konkret werden die Themen Modelle, Skripte, Konfugirationen, generierter Code sowie Firmware-Anbindung über das entwickelte \ac{FOCP} und die Software "FOC Studio".

Nicht behandelt wird die Hardware, also Platine, Messungen, Hardware-Review und Befunde und Inbetriebnahme-Messungen.
% - Dieser Teil: Modelle, Skripte, Konfiguration, generierter Code,
%   Firmware-Anbindung FOCP, PC-Software FOC Studio
% - Hardwareteil (Laborpartner): Platine, Messungen, Hardware-Review
%   (HW-Befunde), Inbetriebnahme-Messungen. Hier nur Verweise darauf.
\offen{Querverweise, Überarbeiten}

\section{Aufbau des Berichts}
\offen{Kurzer Überblick über die Kapitel}
Im Kapitel \ref{ch:ausgangslage} wird zunächst die Ausgangslage, also der vorgefundene Stand der Masterarbeiten von 2023 beschrieben und die Befunde aus der Modellanalyse erläutert.

\ref{ch:vorgehen} stellt die modellbasierte Entwicklung und die verwendeten Werkzeuge vor. Es wird außerdem auf die Konfiguration über Profile und die Versionsverwaltung eingegangen.

Die Überführung des Simulink Modells aus den Masterarbeiten in eine wiederverwendbare Komponentenbibliothek wird in \ref{ch:bibliothek} dargestellt.

Der vorgefundene Regler wies einige Fehler auf. Die Korrektur dieser Befunde behandelt \ref{ch:regler}.

Wie diese Korrekturen durch Model-in-the-Loop-Tests verifiziert wurden ist in \ref{ch:verifikation} beschrieben.

Wie der generierte Quellcode auf die Hardware übertragen wurde und welche Besonderheiten dabei beachtet werden müssen ist in \ref{ch:codegen} festgehalten.

Die Kommunikation mit dem Frequenzumrichter wurde zunächst nur über \ac{CAN} realisiert. Das darunter liegende Protokoll \ac{FOCP} und die aufgetretenen Probleme werden in \ref{ch:kommunikation} diskutiert.

Um den Antrieb gegen Störsignale und kritische Betriebszustände abzusichern, wurden zusätzliche Schutz- und Robustheitsfunktionen ergänzt, darunter eine Begrenzung des Bremsstroms, eine Sollwert-Totzone und eine Überspannungsabschaltung. Diese werden in \ref{ch:schutz} zusammengefasst.

Für die Bedienung des Antriebs am PC wurde die Software FOC Studio entwickelt. \ref{ch:focstudio} dokumentiert ihre Anforderungen, Architektur und Funktionen sowie Test und Auslieferung.

Abschließend zieht \ref{ch:fazit} ein Fazit und benennt die offenen Punkte, die in weiterführenden Arbeiten adressiert werden sollten.
